% Options for packages loaded elsewhere
\PassOptionsToPackage{unicode}{hyperref}
\PassOptionsToPackage{hyphens}{url}
\documentclass[
  ignorenonframetext,
]{beamer}
\newif\ifbibliography
\usepackage{pgfpages}
\setbeamertemplate{caption}[numbered]
\setbeamertemplate{caption label separator}{: }
\setbeamercolor{caption name}{fg=normal text.fg}
\beamertemplatenavigationsymbolsempty
% remove section numbering
\setbeamertemplate{part page}{
  \centering
  \begin{beamercolorbox}[sep=16pt,center]{part title}
    \usebeamerfont{part title}\insertpart\par
  \end{beamercolorbox}
}
\setbeamertemplate{section page}{
  \centering
  \begin{beamercolorbox}[sep=12pt,center]{section title}
    \usebeamerfont{section title}\insertsection\par
  \end{beamercolorbox}
}
\setbeamertemplate{subsection page}{
  \centering
  \begin{beamercolorbox}[sep=8pt,center]{subsection title}
    \usebeamerfont{subsection title}\insertsubsection\par
  \end{beamercolorbox}
}
% Prevent slide breaks in the middle of a paragraph
\widowpenalties 1 10000
\raggedbottom
\AtBeginPart{
  \frame{\partpage}
}
\AtBeginSection{
  \ifbibliography
  \else
    \frame{\sectionpage}
  \fi
}
\AtBeginSubsection{
  \frame{\subsectionpage}
}
\usepackage{iftex}
\ifPDFTeX
  \usepackage[T1]{fontenc}
  \usepackage[utf8]{inputenc}
  \usepackage{textcomp} % provide euro and other symbols
\else % if luatex or xetex
  \usepackage{unicode-math} % this also loads fontspec
  \defaultfontfeatures{Scale=MatchLowercase}
  \defaultfontfeatures[\rmfamily]{Ligatures=TeX,Scale=1}
\fi
\usepackage{lmodern}
\ifPDFTeX\else
  % xetex/luatex font selection
\fi
% Use upquote if available, for straight quotes in verbatim environments
\IfFileExists{upquote.sty}{\usepackage{upquote}}{}
\IfFileExists{microtype.sty}{% use microtype if available
  \usepackage[]{microtype}
  \UseMicrotypeSet[protrusion]{basicmath} % disable protrusion for tt fonts
}{}
\makeatletter
\@ifundefined{KOMAClassName}{% if non-KOMA class
  \IfFileExists{parskip.sty}{%
    \usepackage{parskip}
  }{% else
    \setlength{\parindent}{0pt}
    \setlength{\parskip}{6pt plus 2pt minus 1pt}}
}{% if KOMA class
  \KOMAoptions{parskip=half}}
\makeatother
\setlength{\emergencystretch}{3em} % prevent overfull lines
\providecommand{\tightlist}{%
  \setlength{\itemsep}{0pt}\setlength{\parskip}{0pt}}
% Slide layout defaults for warning-clean scientific decks.
\usepackage{etoolbox}
\IfFileExists{xurl.sty}{\usepackage{xurl}}{}
\IfFileExists{seqsplit.sty}{\usepackage{seqsplit}}{\newcommand{\seqsplit}[1]{#1}}
\protected\def\breakseq#1{\seqsplit{#1}}
\protected\def\breaktt#1{\begingroup\ttfamily\seqsplit{#1}\endgroup}
\setlength{\emergencystretch}{6em}
\tolerance=5000
\hbadness=10000
\hfuzz=1pt
\setlength{\tabcolsep}{2pt}
\AtBeginEnvironment{longtable}{\tiny\renewcommand{\arraystretch}{0.86}\setlength{\tabcolsep}{1pt}}
\AtBeginEnvironment{tabular}{\tiny\renewcommand{\arraystretch}{0.86}\setlength{\tabcolsep}{1pt}}

% Natbib and cross-reference fallbacks — slides don't load natbib
% or cleveref, but combined-PDF manuscript prose may emit these
% commands. The fallback renders citations as a bracketed key list
% and cross-references as detokenized labels so slides don't fail on
% undefined control sequences or raw underscores. \providecommand is
% a no-op if packages load later.
\providecommand{\citep}[1]{[#1]}
\providecommand{\citet}[1]{#1}
\providecommand{\citealp}[1]{#1}
\providecommand{\citeauthor}[1]{#1}
\providecommand{\citeyear}[1]{#1}
\providecommand{\cref}[1]{\texttt{\detokenize{#1}}}
\providecommand{\Cref}[1]{\texttt{\detokenize{#1}}}

\newtheorem{proposition}{Proposition}
\newtheorem{hypothesis}{Hypothesis}
\usepackage{bookmark}
\IfFileExists{xurl.sty}{\usepackage{xurl}}{} % add URL line breaks if available
\urlstyle{same}
% fallback for those not using the hyperref driver hyperxmp:
\makeatletter
\@ifundefined{xmpquote}{\newcommand{\xmpquote}[1]{#1}}{}
\makeatother
\hypersetup{
  hidelinks,
  pdfcreator={LaTeX via pandoc}}

\author{\texorpdfstring{}{}}
\date{}

\begin{document}

\section{Abstract}\label{abstract}

\begin{frame}[fragile,allowframebreaks]{Abstract}
The reproducibility problem in computational research is partly
structural: research artifacts are often distributed across editors,
notebooks, scripts, and manual publication steps without a shared
mechanism that keeps code, data, and manuscript synchronized. Existing
tools address important parts of this lifecycle: workflow managers
orchestrate computation, literate-programming systems render documents,
and data-versioning tools track artifacts. \texttt{template/} combines
those concerns with test and provenance gates by applying Infrastructure
as Code to the research lifecycle. Its Two-Layer Architecture separates
28 infrastructure subdirectories (25 importable Python packages,
\textasciitilde662 modules, validated by \textasciitilde8,583 tests)
from self-contained project workspaces. A YAML-declared pipeline (16
stages; 10 in a default full run) coordinates environment checks, test
execution, analysis, Pandoc/XeLaTeX rendering, SHA-256 hashing with
optional steganographic watermarking, structural validation, and
optional LLM review. The declared quality floors---90\% for project
source and 60\% for infrastructure source---are injected from executable
configuration rather than copied into the manuscript. Documentation
Duality equips repository directories with human-facing
\texttt{README.md} and agent-facing \texttt{AGENTS.md} files;
skill-enabled infrastructure packages additionally expose structured
\texttt{SKILL.md} descriptors so agents can discover supported
capabilities from versioned interfaces.

The generated public exemplar roster
(\breaktt{templates/template\_active\_inference},
\breaktt{templates/template\_autopoiesis},
\breaktt{templates/template\_autoresearch\_project},
\breaktt{templates/template\_autoscientists},
\breaktt{templates/template\_code\_project},
\breaktt{templates/template\_data\_descriptor},
\breaktt{templates/template\_eda\_notebook},
\breaktt{templates/template\_formal},
\breaktt{templates/template\_gold\_refinement},
\breaktt{templates/template\_literature\_meta\_analysis},
\breaktt{templates/template\_madlib},
\breaktt{templates/template\_methods\_paper},
\breaktt{templates/template\_newspaper},
\breaktt{templates/template\_pitch\_deck},
\breaktt{templates/template\_pools\_rules\_tools},
\breaktt{templates/template\_prose\_project},
\breaktt{templates/template\_redacted\_report},
\breaktt{templates/template\_registered\_report},
\breaktt{templates/template\_search\_project},
\breaktt{templates/template\_sia},
\breaktt{templates/template\_storybook},
\breaktt{templates/template\_template},
\breaktt{templates/template\_textbook}) supplies heterogeneous
control-positive layouts, including optimization
(\breaktt{template\_code\_project}, 236 discovered tests), prose
(\breaktt{template\_prose\_project}, 120 discovered tests), and
AutoResearch readiness (\breaktt{template\_autoresearch\_project}, 300
discovered tests). Each declares the same 90\% project-source floor and
consumes the shared pipeline without project-to-project imports. This
manuscript adds a reflexive artifact: \breaktt{template\_template} (138
discovered tests) follows the same analysis and render path while
injecting repository counters and figures from live introspection. The
result is a self-documenting build whose quantitative repository claims
are generated artifacts, whose manuscript claims are checked against an
evidence registry, and whose output can be validated and optionally
watermarked. A versioned comparative matrix situates \texttt{template/}
against peer tools across the capabilities represented in Figure 4. Code
is released under the Apache License 2.0 at
\breaktt{github.com/docxology/template}; the work remains open-ended.
\end{frame}

\end{document}
